\section{Conclusion}

In many domains, people choose to disseminate information across different modalities and formats for better communication to broader audiences. We proposed a unified view of document transformation and evaluation. We showed that LLMs are capable of templatic document generation with minimal supervision, and that a structured, intermediate representation can improve performance, particularly for smaller models. We also introduced a flexible precision-recall framework for automatic evaluation that easily integrates existing evaluation metrics into a unified system and allows for comparison across diverse datasets without additional task specific metric design. Finally, we conducted a human evaluation and showed that annotators prefer the documents generated using the intermediate representation 82\% of the time and that our evaluation framework correlates better with human preference than standard evaluation metrics.% We see this  work as an important step towards the larger goal of generalized document generation. 
%We additionally showed that our framework of evaluation correlates more highly with human preference than standard evaluation metrics.  

% \sjauhar{Consider breaking para here for future work.} The use of templatic views of documents exists in other domains (e.g. creating a cover letter and personal website based on a resume), and in future work we will consider other domains of documents. We position this paper as a step towards the larger goal of many-to-many, multi-modal document template generation. \isabel{I don't like this last sentence, it's not mic drop enough} \sjauhar{Lol, I agree. The last para you had written for the intro (which I commented out) would fit quite well here; consider moving it and adapting into a separate paragraph focused on future work.}


% \sjauhar{Re-iterate the broader impact of this work succintly from the intro.}
